\documentclass[10pt,a4paper]{amsart}
\usepackage{amsmath,amssymb,amsthm,graphicx}
\usepackage{verbatim}
\newenvironment{reference}{\comment}{\endcomment}
%\newenvironment{reference}{}{}
\newenvironment{slogan}{\comment}{\endcomment}
\newenvironment{history}{\comment}{\endcomment}

% For commutative diagrams we use Xy-pic
\usepackage[all]{xy}

% We use 2cell for 2-commutative diagrams.
\xyoption{2cell}
% Theorem environments.
%
\theoremstyle{plain}
\newtheorem{theorem}[subsection]{Theorem}
\newtheorem{proposition}[subsection]{Proposition}
\newtheorem{lemma}[subsection]{Lemma}

\theoremstyle{definition}
\newtheorem{definition}[subsection]{Definition}
\newtheorem{example}[subsection]{Example}
\newtheorem{exercise}[subsection]{Exercise}
\newtheorem{situation}[subsection]{Situation}

\theoremstyle{remark}
\newtheorem{remark}[subsection]{Remark}
\newtheorem{remarks}[subsection]{Remarks}

\numberwithin{equation}{subsection}

% Macros
%
\def\lim{\mathop{\mathrm{lim}}\nolimits}
\def\colim{\mathop{\mathrm{colim}}\nolimits}
\def\Spec{\mathop{\mathrm{Spec}}}
\def\Hom{\mathop{\mathrm{Hom}}\nolimits}
\def\Ext{\mathop{\mathrm{Ext}}\nolimits}
\def\SheafHom{\mathop{\mathcal{H}\!\mathit{om}}\nolimits}
\def\SheafExt{\mathop{\mathcal{E}\!\mathit{xt}}\nolimits}
\def\Sch{\mathit{Sch}}
\def\Mor{\mathop{\mathrm{Mor}}\nolimits}
\def\Ob{\mathop{\mathrm{Ob}}\nolimits}
\def\Sh{\mathop{\mathit{Sh}}\nolimits}
\def\NL{\mathop{N\!L}\nolimits}
\def\CH{\mathop{\mathrm{CH}}\nolimits}
\def\proetale{{pro\text{-}\acute{e}tale}}
\def\etale{{\acute{e}tale}}
\def\QCoh{\mathit{QCoh}}
\def\Ker{\mathop{\mathrm{Ker}}}
\def\Im{\mathop{\mathrm{Im}}}
\def\Coker{\mathop{\mathrm{Coker}}}
\def\Coim{\mathop{\mathrm{Coim}}}

% Boxtimes
%
\DeclareMathSymbol{\boxtimes}{\mathbin}{AMSa}{"02}

%
% Macros for moduli stacks/spaces
%
\def\QCohstack{\mathcal{QC}\!\mathit{oh}}
\def\Cohstack{\mathcal{C}\!\mathit{oh}}
\def\Spacesstack{\mathcal{S}\!\mathit{paces}}
\def\Quotfunctor{\mathrm{Quot}}
\def\Hilbfunctor{\mathrm{Hilb}}
\def\Curvesstack{\mathcal{C}\!\mathit{urves}}
\def\Polarizedstack{\mathcal{P}\!\mathit{olarized}}
\def\Complexesstack{\mathcal{C}\!\mathit{omplexes}}
% \Pic is the operator that assigns to X its picard group, usage \Pic(X)
% \Picardstack_{X/B} denotes the Picard stack of X over B
% \Picardfunctor_{X/B} denotes the Picard functor of X over B
\def\Pic{\mathop{\mathrm{Pic}}\nolimits}
\def\Picardstack{\mathcal{P}\!\mathit{ic}}
\def\Picardfunctor{\mathrm{Pic}}
\def\Deformationcategory{\mathcal{D}\!\mathit{ef}}
\usepackage{cleveref}
\crefname{originalreference}{original reference}{original references}
\Crefname{originalreference}{Original reference}{Original references}
\makeatletter
\newlabel{lemma-cone}{{08C5}{source}{}{sourcepointer-1}{}}
\newlabel{lemma-cone@cref}{{[originalreference][0][]08C5}{[1][0][]source}{}{}{}}
\makeatother
\crefname{originalreference}{original reference}{original references}
\Crefname{originalreference}{Original reference}{Original references}
\makeatletter
\newlabel{lemma-cone-homotopy}{{09QU}{source}{}{sourcepointer-2}{}}
\newlabel{lemma-cone-homotopy@cref}{{[originalreference][0][]09QU}{[1][0][]source}{}{}{}}
\makeatother
\crefname{originalreference}{original reference}{original references}
\Crefname{originalreference}{Original reference}{Original references}
\makeatletter
\newlabel{lemma-cone-rotate-isom}{{09QS}{source}{}{sourcepointer-3}{}}
\newlabel{lemma-cone-rotate-isom@cref}{{[originalreference][0][]09QS}{[1][0][]source}{}{}{}}
\makeatother
\crefname{originalreference}{original reference}{original references}
\Crefname{originalreference}{Original reference}{Original references}
\makeatletter
\newlabel{lemma-restate-axiom-c}{{09QR}{source}{}{sourcepointer-4}{}}
\newlabel{lemma-restate-axiom-c@cref}{{[originalreference][0][]09QR}{[1][0][]source}{}{}{}}
\makeatother
\crefname{originalreference}{original reference}{original references}
\Crefname{originalreference}{Original reference}{Original references}
\makeatletter
\newlabel{situation-ABC}{{09QJ}{source}{}{sourcepointer-5}{}}
\newlabel{situation-ABC@cref}{{[originalreference][0][]09QJ}{[1][0][]source}{}{}{}}
\makeatother
\title{The homotopy category of the differential graded category of Situation ABC is pre-triangulated}
\author{The Stacks Project Authors}
\date{stacks-project:dga.tex:local-snapshot-20261012; GNU FDL 1.2 or later}
\begin{document}
\maketitle
\noindent\textit{Editorial note.} All delivered mathematical text is the author's own, copied from the pinned Stacks Project chapter file sources/stacks-dga/pinned-original-source.tex; no statement or proof step was rewritten, reordered or removed.\par\medskip
\noindent\textit{Editorial note.} The counted proof is the author's complete argument: the TR1 verification from the admissible sequence and cone triangles, the full TR2 argument with the three-row sign-change diagram and both directions, and the TR3 reduction to cone triangles with the completing morphism from lemma-cone are all retained verbatim.\par\medskip
\noindent\textit{Editorial note.} The label lemma-analogue-homotopy-category-pre-triangulated names the counted statement, and the proof environment immediately following it in the source is delivered verbatim as its proof.\par\medskip
\noindent\textit{Editorial note.} The counted statement is carried by the author's own lemma environment in the chapter source.\par\medskip
\noindent\textit{Allowed standard prerequisites (editorial).}\begin{itemize}
\item Distinguished triangles from admissible short exact sequences and cones. Same chapter, Lemmas lemma-cone-homotopy and lemma-cone-rotate-isom. Identifies distinguished triangles with cone triangles up to isomorphism and gives the rotated cone triangle; these two facts are what all three axioms are reduced to.
\item Restated axiom (C) under rotation. Same chapter, Lemma lemma-restate-axiom-c. Closes the converse half of TR2: the rotated cone triangle with sign-changed maps is distinguished.
\end{itemize}\medskip
\ifdefined\hypertarget\hypertarget{sourcepointer-1}{}\fi
\noindent\textit{Source reference.} Original reference 08C5: sources/stacks-dga/pinned-original-source.tex, source lines 3779--3779. Prerequisite pointer invoked at proof line(s) 4569; it names a standard fact of the same chapter that the counted argument uses, and every step of the counted proof built on it is retained verbatim in the delivered body.\par\medskip
\ifdefined\hypertarget\hypertarget{sourcepointer-2}{}\fi
\noindent\textit{Source reference.} Original reference 09QU: sources/stacks-dga/pinned-original-source.tex, source lines 4326--4326. Prerequisite pointer invoked at proof line(s) 4515; it names a standard fact of the same chapter that the counted argument uses, and every step of the counted proof built on it is retained verbatim in the delivered body.\par\medskip
\ifdefined\hypertarget\hypertarget{sourcepointer-3}{}\fi
\noindent\textit{Source reference.} Original reference 09QS: sources/stacks-dga/pinned-original-source.tex, source lines 4167--4167. Prerequisite pointer invoked at proof line(s) 4526; it names a standard fact of the same chapter that the counted argument uses, and every step of the counted proof built on it is retained verbatim in the delivered body.\par\medskip
\ifdefined\hypertarget\hypertarget{sourcepointer-4}{}\fi
\noindent\textit{Source reference.} Original reference 09QR: sources/stacks-dga/pinned-original-source.tex, source lines 4148--4148. Prerequisite pointer invoked at proof line(s) 4556; it names a standard fact of the same chapter that the counted argument uses, and every step of the counted proof built on it is retained verbatim in the delivered body.\par\medskip
\ifdefined\hypertarget\hypertarget{sourcepointer-5}{}\fi
\noindent\textit{Source reference.} Original reference 09QJ: sources/stacks-dga/pinned-original-source.tex, source lines 3762--3762. Prerequisite pointer invoked at proof line(s) 4497; it names a standard fact of the same chapter that the counted argument uses, and every step of the counted proof built on it is retained verbatim in the delivered body.\par\medskip
\begin{lemma}
\label{lemma-analogue-homotopy-category-pre-triangulated}
In Situation \ref{situation-ABC} the homotopy category $K(\mathcal{A})$
with its natural translation functors and distinguished triangles
is a pre-triangulated category.
\end{lemma}

\begin{proof}
We will verify each of TR1, TR2, and TR3. 

\medskip\noindent
Proof of TR1. By definition every triangle isomorphic to a distinguished
one is distinguished. Since
$$
\xymatrix{x\ar[r]^{1_x} & x\ar[r] & 0}
$$
is an admissible short exact sequence, $(x, x, 0, 1_x, 0, 0)$
is a distinguished triangle. Moreover, given a morphism
$\alpha : x \to y$ in $\text{Comp}(\mathcal{A})$, the triangle
given by $(x, y, c(\alpha), \alpha, i, -p)$ is distinguished by
Lemma \ref{lemma-cone-homotopy}.

\medskip\noindent
Proof of TR2. Let $(x,y,z,\alpha,\beta,\gamma)$ be a triangle and
suppose $(y,z,x[1],\beta,\gamma,-\alpha[1])$ is distinguished.
Then there exists an admissible short exact sequence
$0 \to x' \to y' \to z' \to 0$ such that the associated triangle
$(x',y',z',\alpha',\beta',\gamma')$ is isomorphic to
$(y,z,x[1],\beta,\gamma,-\alpha[1])$. After rotating, we conclude
that $(x,y,z,\alpha,\beta,\gamma)$ is isomorphic to
$(z'[-1],x',y', \gamma'[-1], \alpha',\beta')$. By
Lemma \ref{lemma-cone-rotate-isom},
we deduce that $(z'[-1],x',y', \gamma'[-1], \alpha',\beta')$ is
isomorphic to $(z'[-1],x',c(\gamma'[-1]), \gamma'[-1], i, p)$.
Composing the two isomorphisms with sign changes as indicated in
the following diagram:
$$
\xymatrix@C=3pc{
x\ar[r]^{\alpha}\ar[d] &
y\ar[r]^{\beta}\ar[d] &
z\ar[r]^{\gamma}\ar[d] &
x[1]\ar[d] \\
z'[-1]\ar[r]^{-\gamma'[-1]}\ar[d]_{-1_{z'[-1]}} &
x \ar[r]^{\alpha'}\ar@{=}[d] &
y' \ar[r]^{\beta'} \ar[d] &
z'\ar[d]^{-1_{z'}} \\
z'[-1]\ar[r]^{\gamma'[-1]} &
x \ar[r]^{\alpha'} &
c(\gamma'[-1]) \ar[r]^{-p} &
z'
}
$$
We conclude that $(x,y,z,\alpha,\beta,\gamma)$ is distinguished by
Lemma \ref{lemma-cone-homotopy} (2). Conversely, suppose that
$(x,y,z,\alpha,\beta,\gamma)$ is distinguished, so that by
Lemma \ref{lemma-cone-homotopy} (1), it is isomorphic to a
triangle of the form $(x',y', c(\alpha'), \alpha', i, -p)$
for some morphism $\alpha': x' \to y'$ in $\text{Comp}(\mathcal{A})$.
The rotated triangle $(y,z,x[1],\beta,\gamma, -\alpha[1])$ is
isomorphic to the triangle $(y',c(\alpha'), x'[1], i, -p, -\alpha[1])$
which is isomorphic to $(y',c(\alpha'), x'[1], i, p, \alpha[1])$.
By Lemma \ref{lemma-restate-axiom-c}, this triangle is distinguished,
from which it follows that $(y,z,x[1], \beta,\gamma, -\alpha[1])$
is distinguished.

\medskip\noindent
Proof of TR3: Suppose $(x,y,z, \alpha,\beta,\gamma)$ and
$(x',y',z',\alpha',\beta',\gamma')$ are distinguished triangles
of $\text{Comp}(\mathcal{A})$ and let $f: x \to x'$ and
$g: y \to y'$ be morphisms such that
$\alpha' \circ f = g \circ \alpha$. By
Lemma \ref{lemma-cone-homotopy}, we may assume that
$(x,y,z,\alpha,\beta,\gamma)= (x,y,c(\alpha),\alpha, i, -p)$
and $(x',y',z', \alpha',\beta',\gamma')= (x',y',c(\alpha'), \alpha',i',-p')$.
Now apply Lemma \ref{lemma-cone}
and we are done.
\end{proof}

\end{document}
